% ==========================================================
% 060_admissibility_classes.tex
% Forecast Admissibility Surface (FAS)
% ==========================================================

\section{Admissibility Classes}
\label{sec:fas_classes}

The Forecast Admissibility Surface partitions the slice domain into three
qualitatively distinct admissibility classes: \Allowed, \Conditional, and
\Blocked. These classes are categorical governance outcomes, not performance
tiers. Each class encodes a different set of permissions and constraints for
downstream forecasting, evaluation, and readiness control.

\subsection{\Allowed\ Slices}
\label{subsec:fas_allowed}

Slices classified as \Allowed\ exhibit baseline error anatomy consistent with
stable learning and control. Support size is sufficient, spike frequency is low,
and tail behavior is bounded at levels that do not dominate aggregate error or
operational risk. In these slices, baseline instability does not indicate
structural hostility.

For \Allowed\ slices, forecasting workflows may proceed without restriction.
Advanced models may be trained, evaluated, and governed using the full suite of
Forecast Readiness Framework diagnostics. Readiness adjustments and governance
policies may be applied according to standard rules, with the expectation that
learning signals and evaluation metrics will be interpretable and stable.

\subsection{\Conditional\ Slices}
\label{subsec:fas_conditional}

Slices classified as \Conditional\ occupy an intermediate regime. They may
exhibit elevated spike rates, heavier tails, or insufficient observational
support, but do not cross thresholds associated with clear structural
inadmissibility. In many cases, the \Conditional\ classification reflects
uncertainty rather than demonstrated hostility.

Forecasting may still be attempted for \Conditional\ slices, but only under
constrained or conservative rules. Such constraints may include restricted model
classes, limited training horizons, suppressed readiness adjustments, or
heightened governance scrutiny. The purpose of the \Conditional\ class is not to
prohibit forecasting, but to signal that downstream results should be interpreted
with caution.

Importantly, \Conditional\ is not a transient state to be automatically resolved
through additional modeling effort. While some slices may migrate to \Allowed\
as support increases or behavior stabilizes, others may remain perpetually
conditional due to inherent variability.

\subsection{\Blocked\ Slices}
\label{subsec:fas_blocked}

Slices classified as \Blocked\ exhibit baseline error anatomy that is
fundamentally hostile to learning and control. These slices are characterized by
high spike frequency, extreme tail magnitude, or other pathologies that dominate
error behavior and overwhelm stable signal.

For \Blocked\ slices, forecasting and readiness control are explicitly disallowed.
Models trained on such slices are unlikely to generalize, evaluation metrics are
dominated by rare events, and governance signals become erratic or misleading.
Excluding these slices prevents structurally unstable behavior from contaminating
downstream analyses and decision-making.

Blocking is a governance decision, not a judgment of importance. A slice may be
operationally critical yet structurally inadmissible for forecasting. In such
cases, alternative handling—such as deterministic rules, manual planning, or
exogenous overrides—must be employed outside the forecasting pipeline.

\subsection{Class Stability and Interpretation}

Admissibility classes are intended to be stable over time relative to the slice
definition and threshold configuration. While classifications may evolve as new
data becomes available or as slice granularity changes, frequent oscillation is
indicative of a poorly chosen slice domain or overly sensitive thresholds rather
than meaningful structural change.

By separating \Allowed, \Conditional, and \Blocked\ slices, FAS provides a clear
and interpretable governance signal that constrains forecasting behavior without
conflating admissibility with performance, priority, or optimization outcomes.
